% =============================================================================
% style.tex — the palette, the named blocks, and the title page.
%
% Extracted from a course presentation given in 2025. Everything here is the
% author's own; the beamer themes underneath (Madrid, miniframes, circles) ship
% with beamer itself, so there is no third-party licence to carry.
%
% Load after \documentclass and the \usetheme lines. See README.md.
% =============================================================================

% -----------------------------------------------------------------------------
% THE TYPEFACES
%
% Two families with a division of labour, the same arrangement midcentury-olive/
% uses: a text face carries the running matter, a display face carries the
% things you read one line of.
%
%   Red Hat Text     everything you read a paragraph of — body, bullets, block
%                    bodies, tables, the footline.
%   Cormorant        everything you read one line of — the title page, frame
%                    titles, block labels, the section strip.
%
% Both are bundled, both under the SIL Open Font License, both loaded by
% filename from Path.
%
%   fonts/RedHatText/  Red Hat Text, copyright 2019 Red Hat, Inc.
%                      Reserved Font Name "Red Hat" — so it may be
%                      redistributed as it is, which is what happens here, but
%                      a modified copy may not keep the name.
%   fonts/Cormorant/   Cormorant, copyright 2015 Christian Thalmann.
%
% WHY THE SPLIT, RATHER THAN CORMORANT THROUGHOUT
%
% Cormorant is a display face — its own README says so, and its designer drew
% it for large sizes. It is lovely at 21pt and thin at 9pt; a slide's body text
% and its footline are the two places a deck can least afford thin. Red Hat
% Text was drawn for the opposite job, at UI sizes, and holds a projector.
%
% XELATEX ONLY. This is what the fonts cost.
%
% Both are OpenType, OpenType means fontspec, and fontspec means XeLaTeX or
% LuaLaTeX. This template used to build under pdfLaTeX too and no longer does:
% check.sh already defaults to xelatex, but ENGINE=pdflatex fails on the first
% \setsansfont, and the README's "Open in Overleaf" link carries engine=xelatex
% for the same reason.
%
% LOADED BY FILENAME, FROM Path, NOT BY FAMILY NAME.
%
% A family name is resolved by fontconfig against what happens to be installed
% on the machine, and when it is not installed the substitution is SILENT — the
% deck that looked right here is set in something else there. A path that does
% not resolve is a build error. This is the same reason midcentury-olive/
% carries its fonts, and it is why the .otf files are in the repository rather
% than being a thing you are told to install.
%
% CORMORANT MEDIUM, NOT REGULAR. Even used only for display, Regular sits
% lighter than the Red Hat body beside it and the hierarchy inverts — the
% heading reads as the quieter of the two. Medium holds its place.
\usepackage{fontspec}

% The body face goes in the SANS slot and \familydefault follows it, so
% ordinary text — which is most of a deck — needs no markup to be right.
\setsansfont{RedHatText-Regular.otf}[
  Path             = fonts/RedHatText/,
  BoldFont         = RedHatText-Bold.otf,
  ItalicFont       = RedHatText-Italic.otf,
  BoldItalicFont   = RedHatText-BoldItalic.otf,
]
\renewcommand{\familydefault}{\sfdefault}

% The display face goes in the ROMAN slot, so \textrm reaches it — which is
% what \slateTitle and \slateSubTitle already do — and is also given a name of
% its own, for the \setbeamerfont assignments below.
%
% SMALL CAPS COME FROM THE FAMILY AN ELEMENT RESOLVES TO, NOT FROM A DECLARATION
% HERE. Cormorant carries `smcp` in its upright faces, so \textsc works on them
% with nothing else said — verified by removing every SmallCapsFont option and
% rebuilding, which changed not one pixel of the blocks page.
%
% This block used to name CormorantSC explicitly and claim that without it the
% block labels lost their small caps. They did lose them once, but not for that
% reason. Beamer resolves `block title` and `frametitle` through the SANS
% family, which at the time was still TeX Gyre Heros — and Heros has no `smcp`.
% What fixed it is the \setbeamerfont block below, which moves those elements
% onto the display face. See README.md, "small caps that are not small caps".
%
% The real gap, and it is smaller than this note first claimed. Cormorant's
% ITALICS carry no `smcp` and Red Hat Text carries none in any face, both read
% out of the GSUB tables. But a missing `smcp` does not mean \textsc fails:
% rendered and compared, \textsc inside italic Cormorant still SETS small
% capitals — fontspec falls back to the upright face's — so what is actually
% lost is the SLANT, not the small caps. The first version of this note said
% it set ordinary letters, which was inferred from the table and never looked
% at. Nothing in this template asks for small caps in italic, so the fallback
% has never fired here; check it if you write \textsc inside \textit.
\setmainfont{Cormorant-Medium.otf}[
  Path             = fonts/Cormorant/,
  BoldFont         = Cormorant-Bold.otf,
  ItalicFont       = Cormorant-MediumItalic.otf,
  BoldItalicFont   = Cormorant-BoldItalic.otf,
]
\newfontfamily\slateDisplayFont{Cormorant-Medium.otf}[
  Path             = fonts/Cormorant/,
  BoldFont         = Cormorant-Bold.otf,
  ItalicFont       = Cormorant-MediumItalic.otf,
  BoldItalicFont   = Cormorant-BoldItalic.otf,
]

% Which beamer elements are display. Written as full \setbeamerfont
% specifications rather than \addtobeamerfont{}{family=}: the second form
% inherits a series the new family may not have, and fontspec answers a missing
% weight by smearing the outline into a fake one rather than by complaining.
% midcentury-olive/ was bitten by exactly this.
\setbeamerfont{title}             {size=\LARGE,      series=\mdseries, family=\slateDisplayFont}
\setbeamerfont{subtitle}          {size=\large,      series=\mdseries, family=\slateDisplayFont}
\setbeamerfont{frametitle}        {size=\Large,      series=\mdseries, family=\slateDisplayFont}
\setbeamerfont{framesubtitle}     {size=\normalsize, series=\mdseries, family=\slateDisplayFont}
\setbeamerfont{block title}       {size=\normalsize, series=\mdseries, family=\slateDisplayFont}
\setbeamerfont{section in head/foot}{size=\tiny,     series=\mdseries, family=\slateDisplayFont}
\setbeamerfont{title in head/foot}  {size=\tiny,     series=\mdseries, family=\slateDisplayFont}

\usepackage{xcolor}
\usepackage{etoolbox}   % \BeforeBeginEnvironment, used by the block factory
\usepackage{tikz}
\usetikzlibrary{positioning}
\usepackage{makecell}   % multi-line table cells
\usepackage[most]{tcolorbox}   % the SWOT grid at the foot of this file

% -----------------------------------------------------------------------------
% THE RULE UNDER A BLOCK LABEL
%
% Three strokes, not one: a soft wide wash, a hairline over it, and a fainter
% line a hair below. Together they read as a drawn line rather than a printed
% one, which suits a label sitting on a tinted card rather than on a rule.
%
% The marks at the right end are the point of it. Five small parallelograms,
% each a step lighter than the last, standing where the line ends — a hatched
% tail that stops the rule dead rather than letting it run out. They cost
% nothing to draw and they give the block an end.
%
% Adapted from a mockup, and only this part of it. The mockup also came with
% seven new block colours, which were not taken: measured, its Example and
% Comment sat at dE 3.6 from each other — under the 2.3 threshold's neighbour,
% barely a perceptible difference — where the colours already here are 11.9
% apart at their closest. It was the same collision this template had already
% been fixed for.
% -----------------------------------------------------------------------------
\newcommand{\slateblockrule}[1]{%
  \par\nointerlineskip
  % trim=..., and 4pt of it. The strokes have round caps and the hatch marks
  % lean past the end, so the picture measures wider than the line it draws —
  % 4.0pt wider, reported once per block as an Overfull \hbox. Drawing 4pt
  % short of \textwidth costs nothing visible and the log goes quiet.
  \begin{tikzpicture}[x=1pt, y=1pt, baseline=0pt]
    \pgfmathsetmacro{\slrw}{\textwidth-4}
    \draw[#1!15, opacity=0.55, line width=1.7pt, line cap=round] (0,0) -- (\slrw,0);
    \draw[#1!78, line width=0.36pt, line cap=round]              (0,0) -- (\slrw,0);
    \draw[#1!42, opacity=0.36, line width=0.22pt, line cap=round] (0,-0.34) -- (0.97*\slrw,-0.34);
    \foreach \i/\shade in {0/76, 1/68, 2/61, 3/54, 4/48}{%
      \pgfmathsetmacro{\slrx}{\slrw-27.8+5.4*\i}
      \path[fill=#1!\shade, opacity=0.36, rounded corners=0.35pt]
        (\slrx,-0.11) -- ++(4.25,8.2) -- ++(5.1,0) -- ++(-4.25,-8.2) -- cycle;
      \draw[#1!80, opacity=0.38, line width=0.26pt, line cap=round]
        (\slrx+0.57,-0.11) -- ++(3.97,8.2);
    }
    \draw[#1!48, opacity=0.35, line width=0.34pt, line cap=round]
      (\slrw-2.3,-0.34) -- (\slrw,-0.34);
  \end{tikzpicture}%
  \par\nointerlineskip
}

% -----------------------------------------------------------------------------
% HOW A BLOCK IS DRAWN
%
% A label, a rule under it, and a tinted body. No bar, no shadow, no frame.
% Why the bar went is in the note above \blockcolours; this is the geometry.
%
% EVERY SKIP HERE IS EXACT, AND THAT TOOK TWO GOES.
%
% \vskip after \insertblocktitle does nothing useful: the title has not ended
% its paragraph, so the skip forces one, and what lands is a \baselineskip
% rather than the number written. The gap under the label measured 30px at
% 150dpi where 1pt had been asked for. \par then \nointerlineskip then \vskip
% is exact, and \hrule in vertical mode is exact where \rule is not.
%
% tcolorbox brings its own `before skip`, so 4pt above the body rendered as
% 15px. Zeroed here, after which the number written is the number drawn: 1pt
% gives 2px, 4pt gives 8px, 2.5pt gives the 5px this uses.
%
% That accident was also costing a page. The seven-block demo overflowed by
% 87.8pt, and tightening the spacing by hand only ever got it to 46.8pt —
% because the waste was seven stray \baselineskips, not the spacing. Fix the
% skips and the page fits with room over.
%
% coltext=fg IS NOT OPTIONAL. tcolorbox does not inherit beamer's text colour:
% give it colback alone and the body sets in pure black. Sampled from the
% render it came out #000000 where the rest of the deck is #5D677D, and it
% reads as though the block had been deliberately darkened.
% -----------------------------------------------------------------------------
\setbeamertemplate{block begin}{%
  \par\vskip0pt%
  {\usebeamerfont{block title}\usebeamercolor[fg]{block title}\insertblocktitle\par}%
  \nointerlineskip\vskip2pt%
  {\usebeamercolor[fg]{block title}\slateblockrule{fg}}%
  \nointerlineskip\vskip2.5pt%
  \usebeamercolor{block body}%
  \begin{tcolorbox}[colback = bg, colframe = bg, coltext = fg,
                    boxrule = 0pt, arc = 3pt, boxsep = 0pt,
                    left = 6pt, right = 6pt, top = 3pt, bottom = 4pt,
                    before skip = 0pt, after skip = 0pt]%
}
\setbeamertemplate{block end}{\end{tcolorbox}\vskip3pt}

% -----------------------------------------------------------------------------
% THE PALETTE
%
% One colour is chosen; the rest are derived from it by colour-wheel relations,
% so the whole deck re-tints from a single line.
%
%   #bec2cb is a desaturated blue-grey. It is light, which is the point: at
%   full strength it is a background, and at !250 — beamer's syntax for a tint
%   past 100%, i.e. a darkening — it becomes the body text colour. One value
%   doing both ends of the range is what keeps the deck quiet.
%
%   Analogous          neighbours on the wheel. Closest to the theme colour,
%                      so they read as variations rather than contrasts.
%   Tetradic           two complementary pairs. Furthest apart, so these carry
%                      the blocks that must not be confused with each other.
%   Split-complementary  either side of the true complement. A softer contrast
%                      than the complement itself, which at this saturation
%                      would look like a mistake.
%
% All seven sit at the same lightness and saturation. That is deliberate: the
% hue distinguishes them, nothing else, so no block shouts louder than another.
% -----------------------------------------------------------------------------
\definecolor{ThemeColor}{RGB}{190, 194, 203}              % #bec2cb

\definecolor{AnalogousColor-1}{RGB}{193, 190, 203}        % #c1becb
\definecolor{AnalogousColor-2}{RGB}{190, 201, 203}        % #bec9cb

\definecolor{TetradicColor-1}{RGB}{203, 190, 194}         % #cbbec2
\definecolor{TetradicColor-2}{RGB}{203, 199, 190}         % #cbc7be
\definecolor{TetradicColor-3}{RGB}{190, 203, 199}         % #becbc7

\definecolor{SplitComplementaryColor-1}{RGB}{203, 193, 190} % #cbc1be
\definecolor{SplitComplementaryColor-2}{RGB}{201, 203, 190} % #c9cbbe

\setbeamercolor{normal text}{fg=ThemeColor!250}

\setbeamercolor{palette primary}   {bg=ThemeColor,           fg=ThemeColor!12.5!white}
\setbeamercolor{palette secondary} {bg=ThemeColor!87.5!white, fg=ThemeColor!12.5!white}
\setbeamercolor{palette tertiary}  {bg=ThemeColor!75!white,   fg=ThemeColor!12.5!white}
\setbeamercolor{palette quaternary}{bg=ThemeColor!62.5!white, fg=ThemeColor!12.5!white}
\setbeamercolor{structure}    {fg=AnalogousColor-1}
\setbeamercolor{section in toc}{fg=AnalogousColor-1!150}

\setbeamersize{text margin left=10mm, text margin right=10mm}
\setbeamertemplate{navigation symbols}{}

% -----------------------------------------------------------------------------
% NAMED BLOCKS
%
% Beamer gives you three: block, example, alert. This deck wanted seven, each
% naming what it holds — an example, a comment, a proposal, a reminder — and
% each in its own hue.
%
% \newcolouredblock{Environment}{Printed label}{colour} builds one. It uses
% \newtheorem for the numbering and label, and swaps the block colours around
% the environment, because beamer has no per-block colour and \setbeamercolor
% inside a frame leaks into whatever follows.
%
% The original wrote those eight lines out once per block, seven times over.
% The colours are the design; the repetition was not.
% -----------------------------------------------------------------------------
% THE LABEL IS COLOURED, NOT THE BAR BEHIND IT.
%
% Each block used to open with a full-width bar filled at the block's colour
% with the label reversed out in white. That bar was the whole reason unblocked
% text on the same page looked like background: measured on the demo's mixed
% page, a block covered 25.0% of its area in ink against the prose's 8.0% —
% three times the weight, for two paragraphs of equal importance.
%
% The bar was also the only thing that mattered. Taking the body tint away
% instead moved 25.0% to 24.6%; taking the bar away moved it to 7.4%, level
% with the prose. The tint spreads over far more of the page but is too pale to
% register; the bar is small and fully saturated. The shadow, which was the
% first suspect, moved it to 24.9% — nothing.
%
% !300, not the colour itself. These seven are drawn near white, so at full
% strength on their own !25 tint they run 1.48:1 to 1.59:1 — invisible. !250
% leaves three of the seven under 4.5:1 (Proposal 4.18, Problem 4.15, Challenge
% 4.01). !300 puts all seven between 5.98:1 and 8.79:1.
%
% Darkening also makes them easier to tell apart, which is the opposite of what
% you would expect. As pale bars the closest pair — Proposal and Challenge —
% sat at dE 3.3, barely a perceptible difference, so seven "distinguishable"
% block colours were really two or three. At !300 that pair is dE 11.9 and the
% median gap across all pairs goes from 8.1 to 28.6. The seven only became
% seven when the bars came off.
\newcommand{\blockcolours}[1]{%
  \setbeamercolor{block title}{use=example text, fg=#1!300, bg=}%
  \setbeamercolor{block body} {parent=normal text, use=block title example,
                               bg=#1!25}%
}
\newcommand{\blockcoloursreset}{%
  \setbeamercolor{block title}{use=structure, fg=structure.fg!250, bg=}%
  \setbeamercolor{block body} {parent=normal text, use=block title,
                               bg=structure.fg!10!bg}%
}

\newcommand{\newcolouredblock}[3]{%
  \newtheorem{#1}{\textsc{#2}}%
  \BeforeBeginEnvironment{#1}{\blockcolours{#3}}%
  \AfterEndEnvironment{#1}{\blockcoloursreset}%
}

% The seven. Rename them into your own language — the label is the second
% argument and nothing else depends on it.
%
% The environment names carry a suffix. \Example, \Definition and friends are
% already taken: amsmath and beamer's own theorem set define several of them,
% and \newtheorem fails on a name that exists. A suffix is uglier than the bare
% word and cheaper than finding out which words are free in every package a
% future deck might load.
\newcolouredblock{ExampleBox}   {Example}    {ThemeColor}
\newcolouredblock{CommentBox}   {Comment}    {AnalogousColor-1}
\newcolouredblock{ProposalBox}  {Proposal}   {AnalogousColor-2}
\newcolouredblock{ReminderBox}  {Reminder}   {TetradicColor-1}
\newcolouredblock{ProblemBox}   {Problem}    {TetradicColor-2}
\newcolouredblock{ChallengeBox} {Challenge}  {TetradicColor-3}
\newcolouredblock{DefinitionBox}{Definition} {SplitComplementaryColor-1}

% -----------------------------------------------------------------------------
% COVERING THE PAGE WITH A PHOTOGRAPH
%
% \slatecover{file} scales an image to cover the whole slide, whatever its
% aspect ratio, without distorting it. Place the result at the page centre and
% whichever dimension overruns is cropped evenly by the page edge.
%
% \includegraphics[width=\paperwidth] alone does NOT do this. The deck is 16:10
% and most photographs are 16:9; scaled to the paper's width, a 16:9 image
% comes up about a tenth of the page short and leaves a white band across the
% top. Giving both width and height distorts the picture, and adding
% keepaspectratio letterboxes it — the same gap, politely.
%
% So: measure it at full width, and if that is not tall enough, scale by height
% instead.
% -----------------------------------------------------------------------------
\makeatletter
\newsavebox{\slate@coverbox}
\newcommand{\slatecover}[1]{%
  \sbox{\slate@coverbox}{\includegraphics[width=\paperwidth]{#1}}%
  \ifdim\ht\slate@coverbox<\paperheight
    \includegraphics[height=\paperheight]{#1}%
  \else
    \usebox{\slate@coverbox}%
  \fi
}
\makeatother

% -----------------------------------------------------------------------------
% THE WATERMARK
%
% The cover photograph again, on every slide, at 6.7% — far enough down that it
% reads as paper texture rather than as an image. Anything stronger and the
% body text has to fight it.
%
% Comment this out and the deck still works; it is the one part of the design
% that costs nothing to lose.
% -----------------------------------------------------------------------------
\newcommand{\slateWatermark}[1]{%
  \usebackgroundtemplate{%
    \tikz[remember picture, overlay]{%
      \node[opacity=0.067, inner sep=0]
        at (current page.center)
        {\slatecover{#1}};%
    }%
  }%
}

% -----------------------------------------------------------------------------
% THE TITLE PAGE
%
% \slatetitlepage — a full-bleed photograph with the type set straight onto it:
% title and author block in the sky, the date on the ground below the tree
% line, the logo in the top-right corner. Nothing is boxed.
%
% Set the five commands below it first. They are separate commands rather than
% arguments because the title page is the one place a deck wants long strings
% with markup in them, and five arguments to one command is unreadable.
% -----------------------------------------------------------------------------
\newcommand{\slateBackground}{figures/background.png}
\newcommand{\slateTitle}    {\large\textbf{\textrm{Title}}}
\newcommand{\slateSubTitle} {\small\textrm{Subtitle}}
\newcommand{\slateAuthor}   {Author}
\newcommand{\slateAffiliate}{Affiliation}
\newcommand{\slateDate}     {\today}
\newcommand{\slateLogo}     {\includegraphics[width=0.18\paperwidth]{figures/logo.png}}

\newcommand{\slatetitlepage}{%
  \begin{tikzpicture}[remember picture, overlay]

    % opacity=0.55 — the photograph held back far enough that the title box
    % reads as something laid over it rather than as another thing at the same
    % weight. The sky keeps its blue and the maple its red down to about here;
    % below 0.5 the picture stops looking held back and starts looking faded.
    %
    % The usual recipe for this is two layers — the image at 0.65 and a white
    % rectangle at 0.15 over it. Do not bother. The page behind is white, so
    % the veil does nothing the opacity cannot: the two render to within one
    % part in 255 of a single opacity of 0.5525, which was measured, not
    % assumed. One number, one place to change it.
    \node[inner sep=0pt, opacity=0.55] at (current page.center)
      {\slatecover{\slateBackground}};

    % NO PANEL, AND WHY NOT.
    %
    % There used to be a double-stroked, shaded box behind the title here. Its
    % stated job was to lift the title off "a photograph busy enough to swallow
    % a plain panel". Measure the ground it actually sat on and the premise
    % fails: the sky behind it runs at sd 4.2 with no pixel darker than
    % (213,228,245). The panel was armour against a busyness that is not there,
    % and its hard border was the loudest edge on the page.
    %
    % What IS busy is the band of crowns across the middle, from about 0.61 of
    % the page height to 0.86. So the type is placed around that band instead
    % of being armoured against it: this block above it, the date below it.
    % Read the numbers in these comments as worst-case contrast against the
    % darkest single pixel under each line, not against its average.
    %
    % ANY NEW PHOTOGRAPH INVALIDATES ALL OF IT. These offsets are measurements
    % of one picture, not a layout that adapts. Change \slateBackground and the
    % first thing to do is find where its flat ground is.
    %
    % !300 rather than the !250 used elsewhere: with no panel the type is read
    % against the sky directly, and at !250 the subtitle came to 3.58:1, under
    % the 4.5:1 that text this size wants. !300 gives title 5.53:1 and subtitle
    % 5.74:1.
    \node[anchor=north, align=center, text width=0.62\paperwidth,
          color=AnalogousColor-2!300]
      (title) at ([yshift=-0.28\paperheight]current page.north)
      {\LARGE\slateTitle \\[6pt] \small\slateSubTitle};

    % These two stay at !250 — they are still on sky, at 5.28:1 and 5.35:1.
    % \slateAffiliate is the lowest line of the block and clears the crowns by
    % about 0.02\paperheight. That is the whole budget: pushing this block
    % further down puts the affiliation into the trees.
    \node[below=0.55cm] (author)    at (title.south)
      {\color{AnalogousColor-1!250}\slateAuthor};
    \node[below=0.22cm] (affiliate) at (author.south)
      {\color{AnalogousColor-1!250}\small\slateAffiliate};

    % The date is deliberately NOT in that chain. It sits on the field below
    % the tree line, which is the only other flat ground this photograph has:
    % from 0.86 of the page down it runs at sd 8 with nothing darker than 188.
    % Putting it there leaves the trees between the two blocks rather than
    % behind either, and it is why the block above can sit as low as it does.
    %
    % !300 too, but for the opposite reason: the field is darker than the sky,
    % 188 against 232 at worst, and at !250 the date came to 3.32:1. !300 puts
    % it at 5.41:1. In the position it used to occupy — bottom centre, over the
    % maple — it was 1.78:1.
    %
    % So !300 is the LIGHTEST this line can be while it stands on the field,
    % and shrinking it does not relax that — smaller type wants 4.5:1 where
    % \large only wanted 3:1. Anyone who wants it lighter has to move it back
    % up into the sky with the rest of the block, where !250 reaches 5.2:1.
    \node[anchor=south] (date)
      at ([yshift=0.081\paperheight]current page.south)
      {\color{AnalogousColor-1!300}\small\slateDate};

    % Top-right, which is the flattest and palest patch on the page: mean 244,
    % sd 4.4. Bottom centre, where this used to be, is the worst: mean 189,
    % sd 29. Nothing here draws a box — the thin rectangle visible in demo.pdf
    % is inside figures/logo.png, the placeholder, and goes with a real mark.
    \node[anchor=north east, opacity=0.67]
      at ([xshift=-0.045\paperwidth, yshift=-0.055\paperheight]current page.north east)
      {\slateLogo};

  \end{tikzpicture}%
}


% ---- the palette this template mixes from ----------------------------------
%
% The helpful/harmful pair is taken from the two TETRADIC colours, which are
% the furthest apart on the wheel this palette has: greenish against pinkish,
% which is the one contrast a reader will read as good against bad without
% being told. The analogous colours are too close to the theme colour to do
% this job — at this saturation #bec9cb and #bec2cb are the same colour, and
% the two axis headers came out indistinguishable.
%
% Internal/external is the weaker distinction and gets the weaker contrast,
% at a third strength so the row bands stay under the quadrants.
\colorlet{swotHelpful} {TetradicColor-3!67}
\colorlet{swotHarmful} {TetradicColor-1!67}
\colorlet{swotInternal}{ThemeColor!33}
\colorlet{swotExternal}{TetradicColor-2!33}
\colorlet{swotInk}{AnalogousColor-1!250}
\colorlet{S}{swotHelpful!50!swotInternal}
\colorlet{W}{swotHarmful!50!swotInternal}
\colorlet{O}{swotHelpful!50!swotExternal}
\colorlet{T}{swotHarmful!50!swotExternal}

\tcbset{
  swotbox/.style     ={size=normal, boxrule=0pt, colback=#1,
                       watermark text=#1, width=.5\linewidth-5mm},
  swotheader/.style  ={size=small, boxrule=0pt, width=.5\linewidth-5mm,
                       colback=#1, valign=center, halign=center},
  swotfirstcol/.style={swotheader=#1, width=1cm},
}
% =============================================================================
% THE CARD SYSTEM
%
% Brought back from the v5 literature-review deck, where both of these were
% built and settled before the template had either. What arrived is the shape
% they ended at, not the shape they started in — the notes below carry the
% wrong turns, because those are the part worth keeping.
% =============================================================================

% -----------------------------------------------------------------------------
% \slatedivider{label}{headline}{body} — the section-divider card.
%
% Built as one of the deck's own blocks rather than as a new object: the same
% coloured title bar, the same small caps in the display face, so a divider is
% recognisably a member of the family and not a slide from another deck. What
% it adds is a drop shadow, which no other block has -- that is what lifts it
% off the page and says "this is a break", and it is the only place the deck
% uses one.
%
% Three levels inside, each a different job, which is why they are set three
% different ways rather than in three sizes of the same thing:
%
%   label      where you are. Small caps in the display face, in the bar.
%   headline   what this section asks. Display face, one line, dark.
%   body       the requirement itself, verbatim. Text face, grey, so the
%              headline stays the first thing read.
%
% Hyphenation is off in the body. The measure is wide and the paragraph short,
% so a break like "appli-cations" buys nothing and costs the one thing a
% divider page is for, which is being read at a glance. \language=255 rather
% than \hyphenpenalty alone -- see the title page for why the penalty by itself
% does not hold once the font changes inside the argument.
% -----------------------------------------------------------------------------
% \slatedivider{EYEBROW}{Headline}{body} — a whole divider PAGE.
%
% It emits the frame as well as the card: [plain], the card centred between
% two \vspace{\fill}. A deck writes one command where it used to write five
% lines, and the vertical centring stops being something each divider has to
% remember.
%
% THAT IS ALSO THE ONLY REASON IT STILL EXISTS. For a while it was
% \begin{slatecard}...\end{slatecard} and nothing else — one construct wearing
% two names, which is how these two drifted apart in the first place: they were
% two copies of forty lines, with a comment claiming they were shared, and
% flattening one left the other with a filled band and a drop shadow.
%
% Now the two differ in what they ARE. slatecard is a card you put on a page.
% \slatedivider is a page. Anything that needs the card inside a frame of its
% own making — a listing, which cannot be a macro argument at all, or a card
% sharing a page with something else — uses the environment.
\newcommand{\slatedivider}[3]{%
  \begin{frame}[plain]
    \vspace{\fill}
    \begin{slatecard}[slatedividerstyle]{#1}{#2}#3\end{slatecard}
    \vspace{\fill}
  \end{frame}}

% -----------------------------------------------------------------------------
% slatecard — the divider's card as an environment, for content that cannot be
% passed as a macro argument. A verbatim listing is exactly that: lstlisting
% reads its body with catcodes changed, which a \newcommand argument has
% already tokenised by the time the box sees it.
%
% Same shell as \slatedivider, one shared definition, so the two cannot drift
% apart. What differs is only what goes inside.
% -----------------------------------------------------------------------------
% THE BAND AND THE SHADOW ARE GONE; A RULE DOWN THE LEFT EDGE MARKS THE CARD.
%
% The card used to open with a filled ThemeColor band carrying its title, and
% to sit on a drop shadow. Both were doing the same job the seven *Box blocks
% used to do with their bars, and they were measured out for the same reason:
% a filled bar is small but fully saturated, so it dominates a page far beyond
% its area, and the shadow adds ink without adding meaning.
%
% What replaces them is one 3pt rule down the left edge in the darkened theme
% colour, and the title set as a letterspaced eyebrow inside the card rather
% than reversed out of a band above it. The card is still plainly a card: it
% has a fill, a radius and an edge. It just no longer shouts.
%
% sharp corners=west, not all four. The rule reads as an edge only if it is
% straight; rounded off at both ends it becomes a decorative stroke floating
% beside the panel. The right side keeps its radius, which is what the
% reference does and it is right — the marked edge is square, the free edge
% is soft.
\tcbset{slatecardstyle/.style={
    enhanced,
    colback=ThemeColor!10, colframe=ThemeColor!10,
    colbacktitle=ThemeColor!10, coltitle=ThemeColor!300,
    fonttitle=\slateDisplayFont\large,
    % 4mm, not 0. At 0 the eyebrow sat flush against the card's top edge with
    % no air at all, which reads as a clipping error rather than as a design.
    toptitle=4mm, bottomtitle=1.5mm, lefttitle=9mm,
    boxrule=0pt, arc=2.5mm, sharp corners=west,
    borderline west={3pt}{0pt}{ThemeColor!300},
    left=9mm, right=11mm, top=6mm, bottom=8mm,
    width=\linewidth}}

% A DIVIDER IS A FIXED SIZE. A content card may be any height — its content is
% — but a divider is a recurring structural page, and two of them at different
% heights read as a mistake rather than as a rhythm. Measured at 150dpi, the
% divider in this deck came out 340px against 439–527 for the content cards,
% and a second divider with a shorter body would have sat somewhere else again.
%
% `height from`, NOT `height`. Both were built and rendered with a body padded
% until it overran:
%
%   height=58mm            the card stops at 58mm and the TEXT KEEPS GOING,
%                          straight out of the bottom edge and off the slide.
%                          No warning of any kind — not a vbox, not a badness.
%   height from=58mm to    the card grows instead: 340px to 439px at 150dpi.
%                          Still no warning, and correctly so — it grew and it
%                          still fits. Grow it past the frame and check.sh
%                          reports the overflow, which is the whole point.
%
% A hard height is the one thing this template exists to refuse: content
% silently discarded to make a number come out right. Uneven dividers would
% have been the lesser fault.
%
% There is no `minimum height` key in tcolorbox. Asking for one is a pgfkeys
% error that stops the build, which at least fails loudly.
%
% 80mm IS FOUR FIFTHS, AND THE PAGE MAKES THAT EASY: 283.46pt is 100.0mm, so
% the millimetre figure IS the fraction. Do not reach for \paperheight here —
% a macro in this key breaks tcolorbox's parsing of the `to` separator, and the
% error it throws (File ended while scanning use of \pgfkeys@code) points
% nowhere near the cause.
%
% Four fifths rather than five sixths or nine tenths because it is the only one
% that agrees with the page's own margins. Measured off the render, the card's
% left margin is 10.0mm; at 80mm the top and bottom come to 9.0 and 10.9mm, and
% the card reads as a card laid on the page. At 90mm they are 4.5 and 5.5 — the
% card is pressed against the edges and reads as a page with a rule down it.
\tcbset{slatedividerstyle/.style={slatecardstyle, height from=80mm to \textheight}}

% Optional first argument for per-page tcolorbox overrides -- padding, mostly.
% A page carrying a long listing and a footnote needs tighter margins than one
% carrying a paragraph, and the alternative is shrinking the type, which is
% what check.sh exists to prevent.
% \begin{slatecard}[tcb options]{EYEBROW}{Headline}
%
% Two lines of heading, not one. The eyebrow names the KIND of thing the card
% holds — a source, a requirement, a step — and is set small, letterspaced, in
% the sans. The headline is the sentence itself, set large in the display face.
%
% The headline may be empty. \ifblank swallows the whole thing, spacing
% included, so a card that only wants an eyebrow costs nothing and reads
% exactly as it did before this argument existed.
%
% before upper, not a second title. tcolorbox's title is one box with one set
% of paddings; stacking two registers inside it means fighting toptitle and
% bottomtitle for each. Putting the headline at the top of the body instead
% lets the eyebrow keep the title's spacing and the headline take the body's.
\newtcolorbox{slatecard}[3][]{slatecardstyle, #1,
  before lower={},
  fontupper={\language=255\hyphenpenalty=10000\exhyphenpenalty=10000},
  % \textsc, and it works — but only because `fonttitle` above is
  % \slateDisplayFont. The eyebrow is set in Cormorant, which carries `smcp`
  % in its upright faces; accented letters included, checked in the GSUB table
  % and confirmed by rendering a French title.
  %
  % This comment used to say the opposite: that \textsc had to be
  % \MakeUppercase because Red Hat Text has no `smcp`. Red Hat Text has none —
  % that part is true, all four faces — but it is not the face this line
  % reaches. The warning was written for an earlier version where the title
  % resolved through the sans slot, and it outlived it. If you ever drop the
  % \slateDisplayFont from `fonttitle`, the warning comes back: the small caps
  % will quietly become mixed case with the letterspacing still applied, which
  % reads as a near miss rather than a fault.
  title={\addfontfeature{LetterSpace=6}\textsc{#2}},
  before upper={\ifblank{#3}{}{%
    {\slateDisplayFont\LARGE\color{ThemeColor!300}\raggedright #3\par}%
    \vskip2.5mm}\raggedright}}

% -----------------------------------------------------------------------------
% SWOT
%
% \swot{strengths}{weaknesses}{opportunities}{threats} — a 3x3 raster: a blank
% corner, two column headers, and two rows each opening with a rotated label.
%
% The four quadrant colours are MIXED rather than picked, which is what keeps
% the grid readable: every quadrant is half "helpful/harmful" and half
% "internal/external", so the two axes are legible in the colour itself and a
% reader can place a quadrant without reading its label.
%
% The letter watermarked in each quadrant is the tcolorbox colour name. That is
% deliberate — it labels the quadrant without spending a line on a heading.
% -----------------------------------------------------------------------------
% The two ROTATED notes are set at 5pt rather than \tiny, which is 6pt here.
% That one point is what finally cleared the Overfull \hbox this grid had
% reported since it was written — the log is now empty for the first time.
%
% Widening the parbox does not do it. Tried: the warning holds at 6.24pt at
% 2.7cm, 2.9cm and 3.0cm without moving a hair, and by 2.9cm an Overfull \vbox
% appears instead, because the parbox width is the rotated label's height.
% Only making the text narrower works, and there are exactly two ways to do
% that — fewer words or smaller words. midcentury-olive/ took the first and
% says "(of the thing itself)"; this takes the second and keeps the wording.
%
% Worth the point: at 6pt the note wrapped to two lines and sat cramped beside
% the axis name. At 5pt it is one line running parallel to it, which is both
% tidier and what makes the label two lines rather than three.
%
% The HORIZONTAL notes stay at \tiny. They have the full width of a column and
% no reason to shrink.
\newcommand{\swotAxisNote}{\fontsize{5}{6}\selectfont}
\newcommand{\swotAxisHelpful}{Helpful\par\tiny (to achieve the objective)}
\newcommand{\swotAxisHarmful}{Harmful\par\tiny (to achieve the objective)}
\newcommand{\swotAxisInternal}{Internal origin\\\swotAxisNote (attributes of the thing itself)}
\newcommand{\swotAxisExternal}{External origin\\\swotAxisNote (attributes of its environment)}

\newcommand{\swot}[4]{%
  \begin{center}
  \begin{tcbitemize}[raster columns=3, raster rows=3, enhanced, sharp corners,
                     raster valign=center,
                     raster equal height=rows, raster force size=false,
                     raster column skip=0pt, raster row skip=0pt]
    \tcbitem[blankest, width=1cm]
    \tcbitem[swotheader=swotHelpful] {\color{swotInk}\swotAxisHelpful}
    \tcbitem[swotheader=swotHarmful] {\color{swotInk}\swotAxisHarmful}

    % 2.7cm. THIS NUMBER IS A FUNCTION OF THE BODY FONT — re-check it whenever
    % the font changes, because nothing else will tell you.
    %
    % The parbox is rotated, so its width is the length the label has ALONG the
    % row, and the column's own 1cm is how many stacked lines fit across it.
    % The label is three lines: the axis name, then a \tiny parenthetical that
    % wraps to two. If the axis name wraps as well that makes four, the fourth
    % is pushed out of the 1cm column, and it does not overflow into the page —
    % it is simply ABSENT from the render, with nothing in the log naming it.
    %
    % "External origin" is the one that breaks, being a few points longer than
    % "Internal origin", so exactly one of the two labels loses its second line
    % and the page still looks plausible. It has now happened twice:
    %
    %   Computer Modern   2.4cm was enough
    %   TeX Gyre Heros    69.8pt against 2.4cm = 68.3pt   → broke
    %   Red Hat Text      wider still, against 2.5cm      → broke again
    %
    % 2.7cm clears Red Hat Text. Measuring the string in a standalone document
    % is NOT a reliable check — it under-reported by enough to look safe at
    % 2.5cm when it was not. Build the deck and look at the page.
    %
    % WHILE HERE — the Overfull \hbox this grid used to report is gone, and it
    % was never what README.md said it was. It called it a phantom of the
    % raster, on the grounds that narrowing the quadrant boxes never moved it.
    % That evidence was sound and the conclusion was wrong: it was the note in
    % THIS parbox, too wide for its box. Widening the box does not help — see
    % the note above \swotAxisNote — so the note is set a point smaller, and
    % the log is clean.
    \tcbitem[swotfirstcol=swotInternal]
      {\rotatebox[origin=c]{90}{\parbox[t]{2.7cm}{\centering
        \color{swotInk}\swotAxisInternal\par}}}
    \tcbitem[swotbox=S] #1
    \tcbitem[swotbox=W] #2

    \tcbitem[swotfirstcol=swotExternal]
      {\rotatebox[origin=c]{90}{\parbox[b]{2.7cm}{\centering
        \color{swotInk}\swotAxisExternal\par}}}
    \tcbitem[swotbox=O] #3
    \tcbitem[swotbox=T] #4
  \end{tcbitemize}
  \end{center}%
}
